\documentclass[10pt]{article}
\usepackage[margin=0.78in]{geometry}
\usepackage{amsmath,amssymb}
\usepackage[hidelinks]{hyperref}
\usepackage{microtype}
\setlength{\parindent}{0pt}
\setlength{\parskip}{0.45em}

\begin{document}
\begin{center}
  {\Large\bfseries Literature Status: the $1<p<2$ $L_p$ Grothendieck
  Approximation Problem\par}
  \vspace{0.5em}
  {Run \texttt{fa\_banach\_001} \quad 19 July 2026}
\end{center}
\vspace{0.5em}

\textbf{Status:} \texttt{literature\_already\_answered (conditional negative answer)}.
This note records a later literature answer; it is not a new proof produced by
the run.

\section*{Source question}

Khot and Naor define, for an $n\times n$ input matrix $A=(a_{ij})$ with zero
diagonal,
\[
  M_p(A)=\max_{\sum_{k=1}^n |t_k|^p\leq 1}
  \sum_{i=1}^n\sum_{j=1}^n a_{ij}t_it_j.
\]
In Section 5, printed page 26 (PDF page 27), they ask whether, for fixed
$p\in(1,2)$, $M_p(A)$ can be approximated in polynomial time within an
$O(1)$ factor, and state that no hardness-of-approximation result was then
known~\cite{KN}.

\section*{Identification of the answer}

Bhattiprolu, Lee, and Naor prove in Section 8.1, Proposition 8.5(1), (4), and
(5), of the full version of their STOC 2021 paper that the answer is negative
under standard complexity assumptions~\cite{BLN}.  For every fixed $q>2$,
assuming the Small Set Expansion Hypothesis and $P\neq NP$, there is no
polynomial-time constant-factor approximation for the $2\to q$ norm.  Their
Proposition 8.5 transfers this hardness first to the $q^*\to q$ norm and then
to quadratic maximization over the $\ell_{q^*}$ unit ball, explicitly even for
matrices with zero diagonal.

Given any fixed $p\in(1,2)$, take $q=p^*=p/(p-1)>2$.  Then $q^*=p$, so the
last statement is precisely constant-factor inapproximability of $M_p(A)$ in
the zero-diagonal setting of the source question.  The supporting paper also
describes this as hardness for the $\ell_p$ Grothendieck problem that had been
left open.

\section*{Scope limitation}

This is a conditional negative answer.  It rules out the requested algorithm
assuming SSEH and $P\neq NP$; it does not establish an unconditional
complexity separation.  No stronger unconditional claim is made here.

\begin{thebibliography}{9}
\bibitem{KN}
S. Khot and A. Naor,
\emph{Grothendieck-type inequalities in combinatorial optimization},
Commun. Pure Appl. Math. 65 (2012), 992--1035;
\href{https://arxiv.org/abs/1108.2464}{arXiv:1108.2464}.

\bibitem{BLN}
V. Bhattiprolu, E. Lee, and A. Naor,
\emph{A Framework for Quadratic Form Maximization over Convex Sets through
Non-Convex Relaxations}, STOC 2021, 870--881,
\href{https://doi.org/10.1145/3406325.3451128}{doi:10.1145/3406325.3451128};
full draft dated 28 October 2022,
\href{https://www.math.uwaterloo.ca/~vbhattip/quadmax.pdf}{author-hosted PDF}.
\end{thebibliography}

\end{document}
